% Section 4 of the manuscript: the logger.
% Every statement follows the logger's source in kit/logger/ (fixed by the two freezing tags), the format
% specification kit/formats/FORMATS.md, and sections 1.1, 4.4 and 10.3 of the registration. Where the
% registration describes the logger differently, the difference is stated in a footnote.

\section{The logger}
\label{sec:logger}

The logger is a stand-alone program for Windows and macOS. It writes one log per run in a line-based format with a
header, a row for every key press and release of the keys it observes, and a summary row for every conversion
session. The two systems share the format, the key categories and the writing of the rows; they differ in how a
session is recognised, because they offer different routes (Section~\ref{sec:observability}). Table~\ref{tab:rows}
gives the rows and the definitions on each system.

\subsection{Rows}
\label{sec:logger-rows}

\paragraph{Key rows.} For each press and each release of an observed key the logger writes a row with the time on
the wall clock in whole milliseconds, one of fourteen categories and whether the event was a press. The time is read
from the wall clock when the event reaches the logger's keyboard hook (Windows) or event tap (macOS), not taken from the
event's own time stamp, so an adjustment of the system clock moves it. The rows are written for the presses of every
application: the log does not record which application received a key. The categories
keep the kind of a key and discard which key it was: letters, digits, punctuation and symbol keys, the comma and
period keys, Tab, Backspace, Delete, Space, Return, navigation keys, modifier keys, function keys, Escape and other
keys. The comma and period keys form a category of their own because they can mark the boundaries of clauses and
sentences: in Romaji input they type the Japanese comma and full stop. The keys observed differ between the systems (Table~\ref{tab:rows}). On both systems the input-mode keys
produce no row. On macOS rows are written for the letters, the digits and the symbol keys of the main block, Return,
Tab, Space, Backspace, Forward Delete, Escape, the arrow keys and F1 to F12; the modifier keys, the numeric keypad,
Page Up, Page Down, Home, End, the yen and underscore keys of the JIS layout and the function keys above F12 produce
none, and auto-repeated presses are not written.\footnote{The registration describes the macOS logger as writing a row
for every press other than the input-mode keys, F13 to F19 and the modifier keys, and does not mention rows of
releases.
The logger's source observes the narrower set given here and writes a row for each release as well. Apart from the
input-mode keys (the key that the runner pressed before each block to switch the input method on, and the mode keys
of scenario S14), no key outside the observed set was sent in the measurement, and every analysis uses the rows of
presses only.} On Windows rows are
also written for the modifier keys, the numeric keypad, Page Up, Page Down, Home and End; the function keys above F12
produce none, and the driver does not distinguish auto-repeated key-down events. Each synthetic key of the measurement
was held for 60~ms, and the times of releases enter no measure.

\paragraph{Session rows.} For each conversion session the logger writes a row with the duration of the session on
the wall clock (\texttt{dwell\_ms}), the number of arrow-key presses during it (\texttt{nav}) and whether it was
abandoned (\texttt{abandon}). Its time \texttt{t} is the time at which the monitoring loop, which runs every 100~ms,
collected the closed session, up to about 100~ms after the session closed. No row records which key closed a session,
whether a candidate window was shown, or which path recognised the session.

\paragraph{Header.} The first row states the format, the system, the version of the logger, the SHA-256 of the
manifest of the logger's shared source files (the session rules, the key categories and the writing of the rows) and
the row types the log can contain: key rows are written whenever a key is observed, and session rows depend on the
system and its configuration, so that a log without session rows is not read as a log of a person who made no
conversion. The source files of each system's driver are listed in manifests of their own, and the whole source,
drivers included, is fixed by the freezing tags (Section~\ref{sec:deviations}).

\subsection{What the log cannot contain}
\label{sec:logger-content}

No character, reading, candidate string, key name or window name is written. The format is closed: every string in a
row is a member of a fixed enumeration, except the logger's version, which follows a pattern of at most 32 letters,
digits, periods, plus and minus signs, and the SHA-256, which is 64 hexadecimal digits; every number is an integer
within a bounded range; an unknown field, a field given twice and an escape sequence are rejected; and no whitespace
may appear between fields. A validator that is part of the kit
checks this shape on every log, and the analysis and the packing of the deposit refuse a log that fails it. The
validator does not check the meaning of a value: that no content is encoded into a time, a count, the version string or
the SHA-256 of the header is a property of the logger's source, which writes only kinds, durations, counts and times.
The Windows driver reads the class names of windows to decide whether a window belongs to the input method; it does
not keep or write them, and it reads no window title and no accessible name.

\subsection{Sessions on macOS}
\label{sec:logger-macos}

On macOS the logger infers each session from the key presses while a Japanese input source is selected. It learns
the input source from the input-mode keys of the JIS layout once one of them has been seen, and otherwise from polling
the current input source every 100~ms. A session opens at the first press of a letter or digit key. It is closed as
committed by Return, by any key other than a letter, digit, Space, Backspace, arrow or input-mode key, and by a change
to non-Japanese input, at the mode key or at the next press; it is closed as abandoned by Escape. Space, Backspace and the arrow keys keep a session open,
and a press with the Command or Control modifier is a shortcut that leaves the session unchanged. The minus, comma
and period keys, which in Romaji input type a long-vowel mark, a comma and a full stop inside a reading, close a
session, so a reading that contains them is split into several sessions. With the kana layout the same holds for every
kana typed with a symbol key, among them the voiced-sound mark, which the JIS layout places on the @ key; the kana
strings of S12 used only the 39 kana typed with letter, digit, minus, comma and period keys, and the last three of these
close a session as above. A digit that chooses a candidate continues
the session, as does the next letter after a candidate chosen with the mouse, so one row can cover several
conversions. A session is cut off, without a row, at the first press that comes 20~s or more after it opened; that
press starts afresh. \texttt{nav} counts the arrow presses of the session and \texttt{abandon} records that Escape
closed it. These rules are deterministic functions of the presses, so the rows they specify for a known sequence of
presses can be computed (Section~\ref{sec:method-prediction}).

\subsection{Sessions on Windows}
\label{sec:logger-windows}

On Windows a session is a period during which a window of the input method is visible, observed system-wide without
identifying its application. The logger has two paths, chosen from moment to moment. On the first, an accessibility
event reports that a window of the input method has become visible, changed or hidden: a display, or a change while no
session is open, opens a session, and a dismissal closes it. A display while a session is open starts the session
again without a row. This path has no cut-off at 20~s; a window still shown 25~s after the session opened is ended by a
safety valve of the monitoring loop, which writes a row with a duration of about 25~s, and the next change of the
window opens a new session; the valve is part of the logger's main loop (\texttt{kit/logger/src/main.rs}), where a
unit test exercises it, and it wrote no row in the measurement (Section~\ref{sec:results-detection}). On the second path, used until the first window event of a run arrives and whenever none
has arrived for 60~s, the logger follows the cloaking and uncloaking of windows of the input method's class: a session
opens when such a window becomes visible and closes when none is visible, a gap shorter than 150~ms is folded into the
session, and a session is cut off, without a row, 20~s after it opened, whether or not the window is still shown. The monitoring loop
judges whether a candidate window is open from the state that these paths set together with two further signals, the
class of the element that has the keyboard focus and a scan of the visible top-level windows for the class of a
candidate window, and treats any of the three as open. This judgement times the safety valve and closes a session
whose dismissal was missed; on the first path, a raised state that neither further signal confirms is treated as left
over and lowered.\footnote{The registration describes the two paths; the further signals, which are part of the frozen
source, are described here.} On both paths
\texttt{nav} counts the arrow presses and \texttt{abandon} records an Escape press made while a session is open, in any
application, so both are upper bounds for the session. A session that closes and is followed by another before the
monitoring loop collects it is lost.

A row does not record which path wrote it, and which path drove the measured input method was not determined, because
that would have meant reading the logger's session rows before the freeze; the Windows estimates therefore hold for an
unknown mixture of the two paths. Rows of 20 to 25~s can come only from the first path, and rows of about 25~s from its
safety valve.

\subsection{What a researcher can compute from the log}
\label{sec:logger-measures}

Table~\ref{tab:measures} lists, for measures that keystroke studies of writing use, whether each can be computed from
the logger's rows alone on each system, and whether this study estimates its agreement with a reference. The list
covers common measures of the writing literature and the four per-text aggregates that this study defines; it was not
selected by a systematic review. Agreement is estimated for the conversion summaries, the four per-text aggregates
and the timing of the key rows of presses on the software path, and for no other measure.

\begin{table}[t]
\caption{The logger's rows and the definitions on each system}
\label{tab:rows}
\begin{tabular}{L{0.18}L{0.41}L{0.41}}
\toprule
& macOS & Windows \\
\midrule
Key rows & One per press and one per release of an observed key: time (ms), category, press or release & The same \\
Keys observed & Letters, digits, symbol keys of the main block, Return, Tab, Space, Backspace, Forward Delete, Escape, arrows, F1 to F12 & As on macOS, and modifier keys, numeric keypad, Page Up, Page Down, Home, End \\
Keys without a row & Input-mode keys, modifier keys, numeric keypad, Page Up, Page Down, Home, End, JIS yen and underscore keys, keys above F12; auto-repeated presses & Input-mode keys, keys above F12, and keys outside the observed set \\
Session & A segment of kana input inferred from key presses while Japanese input is selected & A period during which a window of the input method is visible (window events, or cloaking) \\
Opens & First letter or digit key & Display of the window (or a change while no session is open); uncloaking of a window of the input method's class \\
Closes & Return or another closing key (committed); Escape (abandoned); a change to non-Japanese input, at the mode key or at the next press after another means of switching & Dismissal of the window; no visible window of the class after a gap of 150~ms \\
Cut-off & At the first press 20~s or more after opening; no row & Window events: none, safety valve at about 25~s with a row; cloaking: 20~s after opening, no row \\
\texttt{dwell\_ms} & From the opening press to the closing press & From the display to the dismissal (window events) or from the uncloaking to the last cloaking (cloaking); to the time of collection when the session is closed by the safety valve or after a missed dismissal \\
\texttt{nav} & Arrow presses during the session & Arrow presses in any application while a session is open (upper bound) \\
\texttt{abandon} & The session was closed by Escape & An Escape press in any application while a session is open (upper bound) \\
\texttt{t} of a session row & Time of collection, up to about 100~ms after the close & The same; on the cloaking path at least 150~ms after the observed end \\
\bottomrule
\end{tabular}
\end{table}

% A long table (main.tex): as a float it is taller than the page of the template.
{\tablebodyfont
\begin{longtable}{L{0.18}L{0.17}L{0.17}L{0.18}L{0.30}}
\caption{Measures of keystroke studies of writing: whether each can be computed from the logger's rows alone, and
whether this study estimates its agreement with a reference}
\label{tab:measures}\\
\toprule
Measure & macOS & Windows & Agreement estimated & Reason when not computable or not estimated \\
\midrule
\endfirsthead
\tablecontinued{5}
\toprule
Measure & macOS & Windows & Agreement estimated & Reason when not computable or not estimated \\
\midrule
\endhead
Inter-key intervals and pauses between presses & Yes & Yes & Timing of the rows of presses on the software path only (Section~\ref{sec:method-checks}) & Physical keyboard and keys outside the observed set not covered \\
Pauses at word and sentence boundaries & Partly & Partly & No & No text; only the comma and period keys are marked, and Space converts rather than separating words \\
Pauses before and after a conversion is committed & Approximate & Approximate & No & The row gives the collection time and the duration, not the time of the commit \\
Key hold times & Without overlap only & Without overlap only & No & Releases pair with presses by category, so overlapping presses of one category cannot be paired; release times not evaluated \\
Revisions (Backspace and Delete) & Counts & Counts & No & No text; whether a deletion edits a reading or committed text is not recorded \\
Bursts of typing & Bounded by pauses only & Bounded by pauses only & No & The conditions on revisions and on the place of insertion need the text and the cursor \\
Keys of one document or application & No & No & No & The application that received a key is not recorded, so key rows of every application are mixed \\
Word counts, production and efficiency ratios & No & No & No & No text and no final product \\
Input mode & No & No & No & The input-mode keys produce no row and the mode is not written \\
Conversions and their duration & Yes, as inferred sessions & As window appearances; one conversion can give several rows & Detection and duration (Section~\ref{sec:method-analysis}) & Different constructs on the two systems \\
Candidate navigation & Yes & Upper bound & Agreement with the arrow presses sent & Counted in any application on Windows \\
Abandonment & Yes & Upper bound & Agreement with the session's outcome & Counted in any application on Windows; in this
study the per-text rate was lower than the reference's (Section~\ref{sec:results-session}) \\
Pointer and mouse actions & No & No & No & Not recorded, so a revision or insertion made with the pointer is not seen \\
Four per-text aggregates of this study & Yes & Yes & Per-text agreement & Defined by this study \\
\bottomrule
\end{longtable}
}
